% Sections 17.1 to 17.5, translated from sv/chapters/17/theory.tex.

\section{Sinusoidal signals and phasors}
\label{c17:sec:fasorer}
A sinusoidal signal and its phasor are related by
\[
  u(t) = \abs{U}\sin(\omega t + \delta) \quad \longleftrightarrow \quad U = \abs{U}\,\angle\,\delta.
\]

\textbf{From sinusoidal signal to phasor:} read the amplitude $\abs{U}$ and the phase angle
$\delta$ directly from the expression for $u(t)$.

\textbf{From phasor to sinusoidal signal:} write $u(t) = \abs{U}\sin(\omega t + \delta)$ with the
phasor $U = \abs{U}\,\angle\,\delta$.

\begin{aside}
\note[Note] All the phasors in a calculation must have the \textbf{same angular frequency}
$\omega$.
\end{aside}

\section{Phasor addition}
\label{c17:sec:fasoraddition}
Sinusoidal signals are most easily added in the phasor plane:
\begin{enumerate}
  \item Convert each sinusoidal signal into a phasor in polar form.
  \item Convert the phasors to rectangular form.
  \item Add the real and imaginary parts.
  \item Convert the sum back to polar form.
  \item Write the sum as a sinusoidal signal.
\end{enumerate}

\section{Multiplication and division with phasors}
\label{c17:sec:multiplikation}
Multiplication in polar form:
\[
  U_1 \cdot U_2 = \abs{U_1}\abs{U_2}\,\angle\,(\delta_1 + \delta_2)
\]
Division in polar form:
\[
  \frac{U_1}{U_2} = \frac{\abs{U_1}}{\abs{U_2}}\,\angle\,(\delta_1 - \delta_2)
\]
These are used in impedance calculations, among other things: $U = Z \cdot I$.

\section{Application: impedance}
\label{c17:sec:impedans}
A voltage $U$ and a current $I$, both complex, give the impedance $Z$:
\[
  Z = \frac{U}{I} = \frac{\abs{U}}{\abs{I}}\,\angle\,(\delta_U - \delta_I)
\]
For the three basic components:
\begin{itemize}
  \item Resistance: $Z_R = R$ (real, phase angle $0$).
  \item Inductance: $Z_L = j\omega L$ (phase angle $+90^{\circ}$).
  \item Capacitance: $Z_C = \dfrac{1}{j\omega C} = -\dfrac{j}{\omega C}$
    (phase angle $-90^{\circ}$).
\end{itemize}

\needspace{12\baselineskip}
\section{Worked examples}
\label{c17:sec:typexempel}

\begin{exempel}[Phasor addition]
Add $u_1(t) = 2\sin(\omega t + 45^{\circ})\,\text{V}$ and
$u_2(t) = 3\sin(\omega t - 60^{\circ})\,\text{V}$.

\noindent\emph{Step 1:} the phasors in polar form:
\[
  U_1 = 2\,\angle\,45^{\circ}, \qquad U_2 = 3\,\angle\,{-60^{\circ}}
\]
\emph{Step 2:} convert to rectangular form:
\begin{gather*}
  U_1 = 2(\cos 45^{\circ} + j\sin 45^{\circ}) = \sqrt{2} + j\sqrt{2} \approx 1.414 + j1.414 \\
  U_2 = 3(\cos(-60^{\circ}) + j\sin(-60^{\circ})) = 3(0.5 - j0.866) = 1.5 - j2.598
\end{gather*}
\emph{Step 3:} add:
\[
  U_{\text{tot}} = (1.414 + 1.5) + j(1.414 - 2.598) = 2.914 - j1.184
\]
\emph{Step 4:} convert to polar form:
\begin{gather*}
  \abs{U_{\text{tot}}} = \sqrt{2.914^2 + 1.184^2} \approx 3.15\,\text{V} \\
  \delta_{\text{tot}} = \arctan\left(\frac{-1.184}{2.914}\right) \approx -22.1^{\circ}
\end{gather*}
\emph{Step 5:} back to the time domain:
\[
  u_{\text{tot}}(t) = 3.15\sin(\omega t - 22.1^{\circ})\,\text{V}
\]
\end{exempel}

\begin{exempel}[Finding the phase from an equation]
An AC voltage is given by $u(t) = 6\sin(80\pi t + \delta)\,\text{V}$. At $t = 15\,\text{ms}$,
$u = 3\,\text{V}$. Find $\delta$.
\begin{gather*}
  6\sin(80\pi \cdot 0.015 + \delta) = 3 \quad \Rightarrow \quad \sin(1.2\pi + \delta) = 0.5 \\
  \delta_1 = \arcsin(0.5) - 1.2\pi \approx 0.5236 - 3.7699 \approx -3.25\,\text{rad} \\
  \delta_2 = \pi - \arcsin(0.5) - 1.2\pi \approx -1.15\,\text{rad}
\end{gather*}
Checking both roots gives $u(0.015) = 3\,\text{V}$.
\end{exempel}

\begin{exempel}[Calculating an impedance]
A series RLC circuit has $R = 10\,\Omega$, $Z_L = j5\,\Omega$ and $Z_C = -j3\,\Omega$ at
$\omega = 1000\,\text{rad/s}$. The total impedance is
\begin{gather*}
  Z = R + Z_L + Z_C = 10 + j5 - j3 = 10 + j2\,\Omega, \\
  \abs{Z} = \sqrt{100 + 4} \approx 10.20\,\Omega, \qquad
  \delta_Z = \arctan\left(\frac{2}{10}\right) \approx 11.3^{\circ}.
\end{gather*}
\end{exempel}
